\begin{afterword}
\summaryhead

The post contrasts two models of belief. Descartes held that we first understand a proposition and then accept or reject it. Spinoza held that understanding a proposition includes accepting it, and that rejecting it takes a further step. If Spinoza is right, distraction should make people remember false statements as true, but not true statements as false.

Daniel Gilbert and colleagues tested this. When people learning made-up words were interrupted, they identified true statements as often as before but false statements much less often (35\% against 55\%). In a second study, people read crime reports in which false statements were marked by color. Those who also had to search for the digit 5 recommended 11.15 years in prison when the false statements made the crime worse and 5.83 years when they excused it; those not searching recommended 7.03 and 6.03 years. The author concludes that we should be careful with unreliable information, especially when doing something else at the same time.

\responsehead

The post reports its two studies accurately. The percentages match the 1990 paper and the prison terms match the 1993 paper's table. It does not claim more than the paper for the non-busy group, whose difference the paper called only marginal. Its closing advice is modest and follows from the crime-report study, whatever explains it: statements marked false changed the judgments of people who were busy.

What the post presents as settled is the explanation. It presents the first study as confirming Spinoza's prediction, and it closes its account of the second with that paper's title, ``You Can't Not Believe Everything You Read.'' The model was disputed when the post appeared. Two years earlier, Hasson, Simmons and Todorov had found Gilbert's asymmetry only for statements that tell you little when false, and concluded that understanding a statement may not require believing it. Later work reported results against the Spinozan model. The reader is not told there is a dispute. That matters more here than in ``Your Strength as a Rationalist,'' where the same research was cited in passing, because this post is the book's account of the experiments. The first study is also small: 35 students, each with two interrupted statements of each kind.

The post's history of philosophy is inaccurate. The post claims that philosophers ``pretty much went along with Descartes.'' Gilbert, Krull and Malone, whom it cites, name Reid and Russell among the philosophers who questioned Descartes, and add that it is modern psychology that kept his view. The footnote on that sentence cites a blog post about policy debates, which does not mention Descartes.

In short: an accurate report of two experiments, whose Spinozan reading is presented as confirmed when it was already disputed, framed by a history of philosophy that the post's own source does not support.
\end{afterword}
